% Layered defense, secure boot, MAC, SBOM
\xchapter{Defense in Depth}

\begin{frame}[plain]
  \vfill\centering
  {\color{tumblue}\Huge\bfseries From Attack to Defense\par}
  \vspace{3mm}
  {\large\color{tumgrey} Secure design, defense in depth, and a security process that never stops}
  \vfill
\end{frame}

\begin{frame}{Defense in Depth: No Single Wall}
  \centering
  \begin{tikzpicture}[font=\scriptsize]
    \def\cx{-2.8}
    % rings, outer -> inner
    \filldraw[fill=tumblue!8,  draw=tumblue!45] (\cx,0) circle(3.40);
    \filldraw[fill=tumblue!16, draw=tumblue!55] (\cx,0) circle(2.72);
    \filldraw[fill=tumblue!26, draw=tumblue!65] (\cx,0) circle(2.04);
    \filldraw[fill=tumblue!40, draw=tumblue!75] (\cx,0) circle(1.36);
    \filldraw[fill=alertred!35, draw=alertred, very thick] (\cx,0) circle(0.80);
    \node[align=center,font=\scriptsize\bfseries] at (\cx,0) {Safety\\ECU};
    \node[font=\scriptsize\bfseries] at (\cx,1.06) {Code};
    \node[font=\scriptsize\bfseries] at (\cx,1.70) {Compiler / OS};
    \node[font=\scriptsize\bfseries] at (\cx,2.38) {OS security};
    \node[font=\scriptsize\bfseries] at (\cx,3.06) {Network};
    \node[align=center,font=\scriptsize,alertred] (atk) at (\cx-3.7,2.7) {remote\\attacker};
    \draw[-Stealth,alertred,very thick] (atk) .. controls (\cx-2.2,2.1) and (\cx-1.4,1.1) .. (\cx-0.5,0.4);
    % legend
    \def\lx{1.5}
    \node[anchor=west] at (\lx,3.0) {\textbf{Network} \,-- gateway, CAN auth (SecOC)};
    \node[anchor=west] at (\lx,2.2) {\textbf{OS} \,-- secure boot, MAC, sandbox};
    \node[anchor=west] at (\lx,1.4) {\textbf{Compiler/OS} \,-- canary, ASLR, NX, CFI};
    \node[anchor=west] at (\lx,0.6) {\textbf{Code} \,-- bounds checks, Rust};
    \node[anchor=west] at (\lx,-0.2) {\textbf{Process} \,-- ISO/SAE 21434, vuln mgmt};
    \node[anchor=west,align=left,alertred] at (\lx,-1.2) {\textbf{Our demo} beat only the inner ring.\\Any outer layer would have stopped it.};
  \end{tikzpicture}
\end{frame}

\begin{frame}{Which Layer Would Have Stopped \emph{Our} Demo?}
  \begin{itemize}
    \item \textbf{Stack canary} -- a secret value just before the return address. Our overflow smashes it, so the program aborts \emph{before} the hijack.
    \item \textbf{ASLR / PIE} -- shuffles addresses every run, so we could never have hardcoded where \texttt{unlock\_ecu()} lives.
    \item \textbf{NX / DEP} -- blocks jumping into injected shellcode (we jumped to existing code instead, which is exactly why return-oriented techniques exist).
    \item \textbf{Bounds-checked copy} (\texttt{strlcpy}, or just checking the length) -- kills the bug at the source. \alert{Cheapest, most reliable fix of all.}
  \end{itemize}
\end{frame}

\begin{frame}{The Security Testing Toolbox}
  \centering
  \begin{tikzpicture}[
      box/.style={draw=tumdark, line width=0.6pt, rounded corners=4pt, fill=tumblue!12,
                  minimum width=3.9cm, minimum height=0.9cm, align=center, font=\small},
      conn/.style={tumdark, line width=0.5pt}]
    \node[circle, fill=tumblue, text=white, minimum size=2.9cm, align=center,
          font=\bfseries] (c) at (0,0) {Security\\Testing\\Activities};
    % left column
    \node[box] (n1) at (-4.5, 2.1) {\textcolor{alertred}{\textbf{1.}}~Source Code Auditing};
    \node[box] (n2) at (-4.5, 0.7) {\textcolor{alertred}{\textbf{2.}}~SW Component Analysis};
    \node[box] (n3) at (-4.5,-0.7) {\textcolor{alertred}{\textbf{3.}}~Vulnerability Scanning};
    \node[box] (n4) at (-4.5,-2.1) {\textcolor{alertred}{\textbf{4.}}~Vulnerability Assessment};
    % right column
    \node[box] (n8) at (4.5, 2.1) {\textcolor{alertred}{\textbf{8.}}~Security Review};
    \node[box] (n7) at (4.5, 0.7) {\textcolor{alertred}{\textbf{7.}}~Security Audit};
    \node[box] (n6) at (4.5,-0.7) {\textcolor{alertred}{\textbf{6.}}~Red Teaming};
    \node[box] (n5) at (4.5,-2.1) {\textcolor{alertred}{\textbf{5.}}~Penetration Testing};
    \foreach \n in {n1,n2,n3,n4,n5,n6,n7,n8}{\draw[conn] (c) -- (\n);}
  \end{tikzpicture}

  \vspace{1mm}
  {\scriptsize\color{tumgrey} A full programme uses many of these. The next slides zoom into two code-level ones: \textbf{SAST} and \textbf{DAST}.}
\end{frame}

\begin{frame}{Where Each Activity Lives: the Secure SDLC}
  \centering
  \begin{tikzpicture}[
      phase/.style={rounded corners=3pt, text=white, font=\scriptsize\bfseries,
                    minimum width=1.8cm, minimum height=0.72cm, text width=1.55cm,
                    align=center, inner sep=2pt},
      act/.style={align=center, font=\scriptsize, text=tumdark, anchor=north, text width=1.92cm}]
    \node[phase, fill=tumblue]     (p1) at (-5,3) {Require\-ments};
    \node[phase, fill=tumblue!88]  (p2) at (-3,3) {Design};
    \node[phase, fill=tumblue!76]  (p3) at (-1,3) {Develop\-ment};
    \node[phase, fill=tumblue!64]  (p4) at (1,3)  {Testing};
    \node[phase, fill=tumblue!52]  (p5) at (3,3)  {Deploy\-ment};
    \node[phase, fill=alertred!80] (p6) at (5,3)  {Mainte\-nance};
    \foreach \a/\b in {p1/p2,p2/p3,p3/p4,p4/p5,p5/p6}{
      \draw[-{Stealth[length=2mm]}, tumgrey, line width=1pt] (\a.east) -- (\b.west);}
    \node[act] at (-5,2.35) {Security\\Strategy \&\\Policy};
    \node[act] at (-3,2.35) {Threat\\Modeling\\[3pt]Designing\\Controls};
    \node[act] at (-1,2.35) {\textbf{Source Code\\Auditing}\\[3pt]SW Component\\Analysis\\[3pt]\textbf{Vulnerability\\Scanning}};
    \node[act] at (1,2.35)  {\textbf{Vulnerability\\Assessment}\\[3pt]\textbf{Penetration\\Testing}\\[3pt]Security\\Review};
    \node[act] at (3,2.35)  {Security\\Audit};
    \node[act] at (5,2.35)  {Continuous\\Pentesting\\[3pt]Red\\Teaming\\[3pt]Continuous\\Monitoring\\[3pt]Incident\\Response};
  \end{tikzpicture}

  \vspace{1mm}
  {\scriptsize\color{tumgrey} \textbf{SAST} fits the \emph{code} activities (Development); \textbf{DAST} and pen-testing fit the \emph{running} ones (Testing). Both shift security \emph{left}, long before a recall.}
\end{frame}

\begin{frame}{Two Ways to Find the Bug: SAST vs.\ DAST}
  \begin{columns}[T]
    \begin{column}{0.5\textwidth}
      \textbf{SAST} -- read the code
      \begin{itemize}
        \item Inspect the source without running it.
        \item A linter such as \texttt{clang-tidy} flags the unchecked \texttt{strcpy} in seconds (e.g.\ its \texttt{clang-analyzer} / \texttt{bugprone} checks).
      \end{itemize}
    \end{column}
    \begin{column}{0.5\textwidth}
      \textbf{DAST} -- break the running thing
      \begin{itemize}
        \item Feed the live program a nasty, oversized input.
        \item \alert{That's exactly what we just did.}
      \end{itemize}
    \end{column}
  \end{columns}
  \vspace{0.3cm}
  \begin{alertblock}{What our live demo was}
    Our demo was DAST in miniature: a UDS-style diagnostic handler that copied attacker input into a fixed buffer with no length check. Input Validation and API Abuse, in a single line of C, caught by either approach.
  \end{alertblock}
  \vspace{0.5mm}
  {\scriptsize\color{tumgrey} \texttt{clang-tidy} and static analysis are covered in depth in \textbf{Davide Cunial}'s lecture (12 November).}
\end{frame}

\begin{frame}{The Part That Never Stops: Vuln Management}
  \begin{itemize}
    \item \textbf{SBOM} -- a full ingredient list of every component in an ECU's software, open-source included.
    \item Cross-checked against CVE feeds non-stop: a new CVE in a library you ship is a hole in your car, whether you wrote that code or not.
    \item Feeds the 21434 loop: triage $\rightarrow$ TARA $\rightarrow$ patch (often OTA) $\rightarrow$ report to the regulator.
    \item A permanent, staffed function inside every OEM now, not a one-off audit before launch.
  \end{itemize}
\end{frame}
